\section{Experiments}
\label{sec:experiments}

We evaluate AnomalyClaw on a 12-domain benchmark spanning industrial, 3D, remote sensing, medical, scene-level, and logical anomaly detection.
All code and evaluation scripts will be released upon publication.

%% ─────────────────────────────────────────────────────
\subsection{Experimental Setup}
\label{sec:setup}

\paragraph{Benchmark.}
We construct a cross-domain evaluation suite (\Cref{tab:benchmark}).
Each domain contributes 120 test images (60 normal, 60 anomalous) drawn from established public datasets, totaling 1{,}440 test items.
Domains are organized sequentially by type: industrial (D1--D4), logical (D5), 3D point cloud (D6), remote sensing (D7), medical (D8--D11), and scene (D12).
VisA (D3) is included specifically for its multi-object categories (PCBs, capsules) and complex texture objects.
Real3D-AD (D6) provides rendered point cloud views testing whether VLMs can detect geometric deformations.
LEVIR-CD+ (D7) introduces satellite imagery change detection, testing cross-modality generalization.

\begin{table}[t]
\centering
\caption{%
    \textbf{Cross-domain benchmark overview.}
    Twelve domains organized by anomaly type.
    Each domain contributes 120 test images (60 normal, 60 anomalous).
}
\label{tab:benchmark}
\small
\setlength{\tabcolsep}{3pt}
\begin{tabular}{@{}clllp{3.0cm}@{}}
\toprule
Code & Domain & Source & Type & Challenge \\
\midrule
D1  & Industrial Mfg.    & MVTec-AD    & Local      & Texture defects \\
D2  & Retail Products    & GoodsAD     & Local      & Packaging damage \\
D3  & Complex Industrial & VisA        & Local+Multi & Multi-object, complex \\
D4  & Infrastructure     & SDNET2018   & Local      & Concrete cracks \\
D5  & Logical Anomaly    & MVTec-LOCO  & Logical    & Rule violations \\
D6  & 3D Product (RGB)   & Real3D-AD   & 3D/Geom.   & Shape deformations \\
D7  & Remote Sensing     & LEVIR-CD+   & Change Det. & Building damage \\
D8  & Dermatology        & DermaMNIST  & Medical    & Subtle morphology \\
D9  & Brain MRI          & BraTS2021   & Medical    & Tissue overlap \\
D10 & Liver CT           & BMAD-Liver  & Medical    & Low-contrast lesions \\
D11 & GI Endoscopy       & HyperKvasir & Medical    & Polyp detection \\
D12 & Road Safety        & BDD100K+RA21 & Scene     & Contextual reasoning \\
\midrule
    & \textbf{Total}           &                   &                      & \textbf{1{,}440 images} \\
\bottomrule
\end{tabular}
\end{table}

\paragraph{Baselines.}
We compare against methods spanning three categories:

\emph{Expert-only methods} (no VLM):
\begin{enumerate}[leftmargin=*,itemsep=1pt]
    \item \textbf{CLIP-ZeroShot}: Zero-shot anomaly classification using CLIP text-image similarity with prompt templates (``a photo of a normal/defective [object]'').
    \item \textbf{WinCLIP}~\citep{jeong2023winclip}: Window-level CLIP similarity with multi-scale aggregation.
    \item \textbf{PatchCore Expert}: Our patch-level expert in isolation (DINOv2 features, top-1\% aggregation, $k{=}4$ retrieved references).
\end{enumerate}

\emph{Single-pass VLM methods} (one VLM call, no debate):
\begin{enumerate}[leftmargin=*,itemsep=1pt,start=4]
    \item \textbf{VLM-Direct}: VLM receives query image only, no references or expert evidence.
    \item \textbf{Retrieval+VLM}: DINOv2 retrieval followed by VLM comparison between query and references.
    \item \textbf{Expert-Informed VLM}: Retrieval + expert evidence as text context, single-pass VLM judgment (AnomalyClaw without debate).
\end{enumerate}

\emph{Multi-round/agent methods}:
\begin{enumerate}[leftmargin=*,itemsep=1pt,start=7]
    \item \textbf{AgentIAD}~\citep{zhang2025agentiad}: Tool-augmented single agent with iterative reasoning.
    \item \textbf{Symmetric Debate}: Two VLM agents with identical roles debating anomaly presence (no Proposer/Advocate asymmetry, no expert grounding).
\end{enumerate}

\paragraph{Implementation details.}
Unless otherwise noted, experiments use \tbd{[primary VLM]} as the debate backbone with temperature~0 and structured JSON output.
The patch expert uses DINOv2 ViT-S/14 at native $518{\times}518$ resolution.
Default debate depth is $D_{\max}=2$.
We report AUROC as the primary metric.
Confidence intervals are 95\% bootstrap CIs over 10{,}000 resamples.


%% ─────────────────────────────────────────────────────
\subsection{Main Results}
\label{sec:main_results}

\Cref{tab:main_results} presents per-domain AUROC for all methods.

\begin{table}[t]
\centering
\caption{%
    \textbf{Per-domain AUROC on the 12-domain benchmark.}
    \colorbox{bestcolor}{Green} = best; \colorbox{secondcolor}{orange} = second best.
    AnomalyClaw results use \tbd{[primary VLM]} with debate depth $D{=}2$.
}
\label{tab:main_results}
\scriptsize
\setlength{\tabcolsep}{2pt}
\begin{tabular}{@{}lcccccccccccc|c@{}}
\toprule
Method & D1 & D2 & D3 & D4 & D5 & D6 & D7 & D8 & D9 & D10 & D11 & D12 & Macro \\
\midrule
\multicolumn{12}{@{}l}{\emph{Expert-only (no VLM)}} \\
CLIP-ZeroShot
    & 	bd{.XX} & 	bd{.XX} & 	bd{.XX} & 	bd{.XX} & 	bd{.XX} & 	bd{.XX}
    & 	bd{.XX} & 	bd{.XX} & 	bd{.XX} & 	bd{.XX} & 	bd{.XX} & 	bd{.XX} & 	bd{.XXX} \\
WinCLIP
    & 	bd{.XX} & 	bd{.XX} & 	bd{.XX} & 	bd{.XX} & 	bd{.XX} & 	bd{.XX}
    & 	bd{.XX} & 	bd{.XX} & 	bd{.XX} & 	bd{.XX} & 	bd{.XX} & 	bd{.XX} & 	bd{.XXX} \\
PatchCore Expert
    & 	bd{.XX} & 	bd{.XX} & 	bd{.XX} & 	bd{.XX} & 	bd{.XX} & 	bd{.XX}
    & 	bd{.XX} & 	bd{.XX} & 	bd{.XX} & 	bd{.XX} & 	bd{.XX} & 	bd{.XX} & 	bd{.XXX} \\
\midrule
\multicolumn{12}{@{}l}{\emph{Single-pass VLM}} \\
VLM-Direct
    & 	bd{.XX} & 	bd{.XX} & 	bd{.XX} & 	bd{.XX} & 	bd{.XX} & 	bd{.XX}
    & 	bd{.XX} & 	bd{.XX} & 	bd{.XX} & 	bd{.XX} & 	bd{.XX} & 	bd{.XX} & 	bd{.XXX} \\
Retrieval+VLM
    & 	bd{.XX} & 	bd{.XX} & 	bd{.XX} & 	bd{.XX} & 	bd{.XX} & 	bd{.XX}
    & 	bd{.XX} & 	bd{.XX} & 	bd{.XX} & 	bd{.XX} & 	bd{.XX} & 	bd{.XX} & 	bd{.XXX} \\
Expert-Informed VLM
    & 	bd{.XX} & 	bd{.XX} & 	bd{.XX} & 	bd{.XX} & 	bd{.XX} & 	bd{.XX}
    & 	bd{.XX} & 	bd{.XX} & 	bd{.XX} & 	bd{.XX} & 	bd{.XX} & 	bd{.XX} & 	bd{.XXX} \\
\midrule
\multicolumn{12}{@{}l}{\emph{Multi-round / agent}} \\
AgentIAD
    & 	bd{.XX} & 	bd{.XX} & 	bd{.XX} & 	bd{.XX} & 	bd{.XX} & 	bd{.XX}
    & 	bd{.XX} & 	bd{.XX} & 	bd{.XX} & 	bd{.XX} & 	bd{.XX} & 	bd{.XX} & 	bd{.XXX} \\
Symmetric Debate
    & 	bd{.XX} & 	bd{.XX} & 	bd{.XX} & 	bd{.XX} & 	bd{.XX} & 	bd{.XX}
    & 	bd{.XX} & 	bd{.XX} & 	bd{.XX} & 	bd{.XX} & 	bd{.XX} & 	bd{.XX} & 	bd{.XXX} \\
\rowcolor{gray!10}
\textbf{AnomalyClaw (Ours)}
    & 	bd{.XX} & 	bd{.XX} & 	bd{.XX} & 	bd{.XX} & 	bd{.XX} & 	bd{.XX}
    & 	bd{.XX} & 	bd{.XX} & 	bd{.XX} & 	bd{.XX} & 	bd{.XX} & 	bd{.XX} & 	bd{.XXX} \\
\bottomrule
\end{tabular}
\end{table}

\paragraph{Key observations.}
\tbd{[Analysis of main results to be written after experiments complete.
Expected observations:
(1) AnomalyClaw should outperform all baselines on macro AUROC.
(2) The gap over single-pass VLM methods should be larger than in the previous pipeline design, since debate catches hallucinations.
(3) Expert-only baselines (CLIP-ZeroShot, WinCLIP) should be substantially weaker, establishing a clear progression.
(4) Symmetric debate should underperform AnomalyClaw, validating the asymmetric design.
(5) Largest gains expected on D5 (logical), D3 (multi-object), and medical domains where self-verification matters most.]}


%% ─────────────────────────────────────────────────────
\subsection{Ablation Studies}
\label{sec:ablation}

\paragraph{Ablation 1: Does adversarial debate help?}
\Cref{tab:ablation_debate} isolates the contribution of the debate mechanism by comparing AnomalyClaw variants with identical expert evidence but different reasoning strategies.

\begin{table}[t]
\centering
\caption{%
    \textbf{Debate ablation.}
    Adversarial debate with expert grounding outperforms all single-pass and symmetric alternatives.
}
\label{tab:ablation_debate}
\small
\begin{tabular}{@{}llcc@{}}
\toprule
Reasoning Strategy & Expert Evidence & Macro AUROC & VLM Calls/Query \\
\midrule
Single-pass VLM             & None           & \tbd{X.XXX} & 1 \\
Single-pass VLM             & Patch + Retrieval & \tbd{X.XXX} & 1 \\
Symmetric debate (no expert) & None           & \tbd{X.XXX} & 4 \\
Symmetric debate            & Patch + Retrieval & \tbd{X.XXX} & 4 \\
AnomalyClaw (asymmetric)     & Patch + Retrieval & \tbd{X.XXX} & 2--4 \\
\bottomrule
\end{tabular}
\end{table}

\paragraph{Ablation 2: Which experts matter?}
\Cref{tab:ablation_experts} tests the contribution of different experts within the AnomalyClaw framework.

\begin{table}[t]
\centering
\caption{%
    \textbf{Expert pool ablation.}
    Patch-level evidence provides the largest single improvement; combining multiple experts yields additional gains.
}
\label{tab:ablation_experts}
\small
\begin{tabular}{@{}lc@{}}
\toprule
Expert Configuration & Macro AUROC \\
\midrule
No experts (debate only)                & \tbd{X.XXX} \\
Retrieval expert only                    & \tbd{X.XXX} \\
Patch expert only                        & \tbd{X.XXX} \\
Patch + Retrieval                        & \tbd{X.XXX} \\
Patch + Retrieval + Texture              & \tbd{X.XXX} \\
\bottomrule
\end{tabular}
\end{table}

\paragraph{Ablation 3: Debate depth.}
\Cref{tab:ablation_depth} varies the maximum debate depth $D_{\max}$.

\begin{table}[t]
\centering
\caption{%
    \textbf{Effect of debate depth.}
    Performance improves from $D{=}1$ to $D{=}2$ and plateaus at $D{=}3$.
}
\label{tab:ablation_depth}
\small
\begin{tabular}{@{}ccc@{}}
\toprule
$D_{\max}$ & Macro AUROC & Avg.\ VLM Calls/Query \\
\midrule
1 & \tbd{X.XXX} & 2.0 \\
2 & \tbd{X.XXX} & \tbd{X.X} \\
3 & \tbd{X.XXX} & \tbd{X.X} \\
4 & \tbd{X.XXX} & \tbd{X.X} \\
\bottomrule
\end{tabular}
\end{table}


%% ─────────────────────────────────────────────────────
\subsection{Multi-VLM Generalization}
\label{sec:multi_vlm}

A key claim of AnomalyClaw is VLM-agnostic design.
\Cref{tab:multi_vlm} evaluates the framework across multiple VLM backends.

\begin{table}[t]
\centering
\caption{%
    \textbf{Cross-VLM generalization.}
    AnomalyClaw consistently improves over single-pass baselines regardless of the underlying VLM.
}
\label{tab:multi_vlm}
\small
\begin{tabular}{@{}lcc@{}}
\toprule
VLM Backend & Single-Pass AUROC & AnomalyClaw AUROC \\
\midrule
\tbd{VLM-A}          & \tbd{X.XXX} & \tbd{X.XXX} \\
\tbd{VLM-B}          & \tbd{X.XXX} & \tbd{X.XXX} \\
\tbd{VLM-C (open)}   & \tbd{X.XXX} & \tbd{X.XXX} \\
\bottomrule
\end{tabular}
\end{table}


%% ─────────────────────────────────────────────────────
\subsection{Analysis}
\label{sec:analysis}

\paragraph{Cost--accuracy tradeoff.}
The autonomous controller enables a natural cost--accuracy tradeoff through debate depth control.
\tbd{[Pareto plot comparing: PatchCore Expert (0 VLM calls), single-pass VLM (1 call), AnomalyClaw D=1 (2 calls), AnomalyClaw D=2 (2--4 calls), symmetric debate (4 calls).
AnomalyClaw should Pareto-dominate symmetric debate (better accuracy, similar or lower cost).]}

\paragraph{When does debate help most?}
\tbd{[Per-domain analysis showing debate gain over single-pass.
Expected: largest gains on D5 (logical---debate catches compositional errors), D3 (multi-object---debate resolves ambiguous multi-region claims), and medical domains (debate catches hallucinated defects on normal tissue).
Smallest gains on D10 (road safety---semantically obvious anomalies).]}

\paragraph{Autonomous controller behavior.}
\tbd{[Statistics on expert selection frequency and debate depth per domain.
Expected: D10 mostly 1-round (easy), medical domains mostly 2-round (ambiguous), D5 frequently triggers all experts (complex).]}

\paragraph{Failure analysis.}
\tbd{[Identify worst-performing domain (likely D8 Liver CT) and analyze error modes.
Focus on: Does debate reduce false positives? Does it help with false negatives? What types of anomalies resist adversarial verification?]}
